% ============================================================
%  D flip-flop (master-slave, edge-triggered) - logic diagram.
%  Two D latches in series + one inverter on the clock:
%     master latch  enable = CLK-bar  (captures D while CLK = 0)
%     slave  latch  enable = CLK       (passes it out while CLK = 1)
%  Result: Q takes the value of D on the RISING edge of CLK,
%  then ignores D until the next edge.
%  Compile:  pdflatex circuit.tex
% ============================================================
\documentclass[border=18pt]{standalone}
\usepackage{circuitikz}
\usepackage{amsmath}
\usetikzlibrary{calc}

\tikzset{dl/.style={thick, rounded corners=2pt, fill=blue!5}}

\begin{document}
\begin{circuitikz}[line width=1.1pt, font=\sffamily]
  \ctikzset{logic ports=american, logic ports/thickness=1.1}

  % ---------- master and slave D-latch blocks ----------
  \draw[dl] (2.8,1.4) rectangle (6.0,4.4);
  \node[font=\bfseries] at (4.4,3.5) {D latch};
  \node[font=\footnotesize] at (4.4,2.9) {(master)};
  \draw[dl] (9.4,1.4) rectangle (12.6,4.4);
  \node[font=\bfseries] at (11.0,3.5) {D latch};
  \node[font=\footnotesize] at (11.0,2.9) {(slave)};

  % port stubs + labels for the two latch boxes
  \foreach \bx/\by in {2.8/0, 9.4/0}{
    \node[anchor=west,font=\footnotesize] at (\bx+0.1,3.8) {D};
    \node[anchor=west,font=\footnotesize] at (\bx+0.1,2.0) {E};
  }
  \node[anchor=east,font=\footnotesize] at (5.9,3.8) {Q};
  \node[anchor=east,font=\footnotesize] at (12.5,3.8) {Q};
  \node[anchor=east,font=\footnotesize] at (12.5,2.0) {$\overline{\mathrm Q}$};

  % ---------- D input ----------
  \draw (1.0,3.8) node[ocirc]{} node[left]{$D$} -- (2.8,3.8);

  % ---------- master Q -> slave D ----------
  \draw (6.0,3.8) -- (9.4,3.8);
  \node[anchor=south,font=\footnotesize] at (7.7,3.85) {$Q_{\mathrm{master}}$};

  % ---------- slave outputs Q, Qbar ----------
  \draw (12.6,3.8) -- (14.4,3.8) node[ocirc]{} node[right]{$Q$};
  \draw (12.6,2.0) -- (14.4,2.0) node[ocirc]{} node[right]{$\overline{Q}$};

  % ---------- clock + inverter ----------
  \node[not port] (inv) at (2.3,-1.0) {};
  \node[below=1pt,font=\scriptsize] at (inv.south) {clock inverter};
  \draw (-0.4,-1.0) node[ocirc]{} node[left]{CLK} -- (0.2,-1.0) coordinate(clk);
  \fill (clk) circle (2pt);
  \draw (clk) -- (inv.in);
  % CLK-bar -> master E
  \draw (inv.out) -- (2.6,-1.0) -- (2.6,2.0) -- (2.8,2.0);
  \node[anchor=west,font=\footnotesize] at (2.65,0.6) {$\overline{\mathrm{CLK}}$};
  % CLK -> slave E
  \draw (clk) -- (0.2,0.6) -- (9.0,0.6) -- (9.0,2.0) -- (9.4,2.0);
  \node[anchor=south,font=\footnotesize] at (6.5,0.62) {CLK};

\end{circuitikz}
\end{document}
